\section{Introduction}

\subsection{Background}

The Research Centre for Data Intelligence conducts applied research into, among
other themes, smart assistants and human-computer interaction. This project
was carried out in the context of the centre's ``Measuring at Home with AI''
research direction. The long-term motivation for that research is to make
exercise monitoring more accessible to children with juvenile dermatomyositis
(JDM). These children currently perform exercises related to the Childhood
Myositis Assessment Scale (CMAS) under the supervision of specialists at the
Wilhelmina Children's Hospital (WKZ). Repeated travel to the hospital is
burdensome for children and their families, which creates an opportunity for
technologies that can support monitoring at home.

Automated home monitoring must do more than record a video. It must detect a
person, estimate body movement from visual input, distinguish between
activities, and communicate its interpretation clearly to the user. It must
also operate under practical conditions, such as changing distances,
incomplete views of the body, and limited computing resources. The stakeholders
for this feasibility study are the Research Centre for Data Intelligence, the
project client and researcher, and the children and families who represent the
future users of a home-monitoring system. The project requirements were
consolidated in a software requirements specification covering perception,
robot interaction, activity recognition, feedback, session management, privacy,
accuracy, reliability, usability, and safety. The current project addresses the
technical feasibility of such a system; it does not develop or validate a
medical assessment instrument.

\subsection{Opportunity}

A smartphone or tablet could provide a camera for movement monitoring, but a
robot offers an additional interaction layer. The Reachy Mini can follow a
person with its head and body, process camera input on its own computer, and
respond through head and antenna movements. These behaviours make it possible
to present the system as an exercise buddy rather than only as a measurement
device. Such expressive feedback may help make an interaction more
understandable and engaging, particularly for children exercising without a
specialist in the room. Engagement, adherence, and clinical acceptability are
important future questions, but they are outside the evaluation of this
prototype.

\subsection{Goal and scope}

The goal of this project is to investigate whether the Reachy Mini can perform
person detection and activity monitoring, including expressive feedback, under
the computational constraints of the wireless robot. The resulting artefact is
a working prototype that runs in simulation and on the physical robot. It
contains a shared camera-to-pose-to-classifier pipeline, temporal feature
extraction, activity-recognition models, and robot applications that translate
recognized actions into antenna responses. Candidate vision models were
benchmarked under the robot's CPU constraints, leading to the selection of
MoveNet int8 for on-robot pose estimation. Model-training and export tools were
also developed so that the robot can perform inference with NumPy-only model
files rather than requiring the full training environment. Recognition models
were evaluated on public skeleton datasets with subject or footage-group
separation, as well as on robot recordings and an external dataset. Users can
also record personalized movements and create their own datasets for future
recognition models.

The prototype recognizes six classes: none, wave, pushup, squat, clapping,
and jumping jacks. Evaluation focuses on technical feasibility using
simulation, public skeleton datasets, recorded robot footage, an external test
dataset, and different camera framings. The robot's camera should ideally
capture the whole person, because waist-up framing removes information that is
important for some activities.
The prototype is not a CMAS scorer, is not clinically validated, and should
not be used as a clinical decision-support system.

The central research question is:

\begin{quote}
Can the Reachy Mini serve as an engaging, CPU-only platform for home activity
monitoring, and what are the technical and practical limits of the current
prototype?
\end{quote}

The remainder of this article first discusses the theoretical foundations and
related solution space. It then describes the prototype and its evaluation
method, presents the results, discusses their implications and limitations,
and concludes with directions for future development.
